\section{Learning-based generation}
\label{sec:coral-island-cGAN-training}

In this section, we introduce the use of a cGAN, specifically the pix2pix model, to enhance the island generation process by increasing the variety and flexibility of terrains. While the initial procedural algorithm generates diverse island examples (convex and concave outlines), cGAN provides additional flexibility in generating more complex terrain without the rigid constraints of the procedural algorithm that stem from our initial assumptions based on coral reef formation theory.

Our model is trained using a batch size of 1, and optimised using the Adam optimiser with a learning rate of 0.0002 and momentum parameters $\beta_1 = 0.5$ and $\beta_2 = 0.999$. The loss function combined a conditional adversarial loss and an L1 loss, where the L1 loss was weighted by a factor $\lambda = 100$. The generator follows a U-Net architecture with encoder-decoder layers and skip connections, while the discriminator was based on a PatchGAN, classifying local image patches. These parameters are directly adopted from the original pix2pix paper \cite{Isola2017}. We use a dataset of 1~000 single islands, augmented with the various methods seen in \cref{sec:coral-island-data-augmentation} to obtain about 10~000 samples that are all seen once during an epoch (representing about 30 minutes of training).  

During inference after a single epoch, our model still outputs grid artefacts, visible in \cref{fig:coral-island-first-epoch}. After the second epoch, visual artefacts are already rare. Illustrations displayed in the remainder of this chapter are results of models trained over less than ten epochs. The training time is relatively short compared to typical deep-learning terrain generation pipelines, largely thanks to the synthetic nature and consistency of our dataset.
% This training time seems long but, to the best of our knowledge, is the shortest training time for terrain generation with a learning-based approach in the literature, mostly due to the synthetic nature of our dataset.

\begin{figure}%[t]
    \autofitgraphics[]{epoch-1-output.png, epoch-2-output.png, epoch-3-output.png}
    \autofitgraphics[]{epoch-1-render.png, epoch-2-render.png, epoch-3-render.png}
    \autofitcaptions[]{Epoch 1 \linebreak ($\sim$10~000 samples), Epoch 1.5 \linebreak ($\sim$15~000 samples), Epoch 2 \linebreak ($\sim$20~000 samples)}
    \caption[Outputs of the cGAN after training on 10~000, 15~000 and 20~000 samples]{After the first epoch (left), grid artefacts similar to a low-resolution image are visible, but are being corrected during a second epoch (middle) where erosion patterns appear in the height fields (right).}
    \label{fig:coral-island-first-epoch}
\end{figure}

Once the model is trained, the user colours a $256\times256$ RGB image to sketch the different regions. Each region is given a specific colour based on the HSV encoding used for the dataset generation. The examples presented in this paper use red for abysses, blue for reefs, cyan for lagoons, green for beaches, and yellow for mountains. The subsidence factor presented in \cref{sec:coral-island-subsidence} is controllable through the luminosity of the input image.

The Python script for the initial island dataset generation is unoptimised and takes about 2.5 s per island of size $256 \times 256$ as we do not perform parallelisation here. Implementing an optimised C++ version of the initial generation process reduces this execution time to 50 ms per generation.

On the other hand, the inference time of our deep learning model for a single input image of dimension $256 \times 256$ remains constant regardless of scene complexity. Using the NVIDIA GeForce GTX 1650 Ti GPU with Python 3.10 and PyTorch version 2.5.1+cu121, the inference time measured is 5 ms (std 1.1 ms).

Once trained, the cGAN model generates a height field from a given label map. The output is a complete 2D height map representing the terrain's elevation at each point. Although the cGAN's internal logic is not easily interpretable, the label map remains available after inference and retains semantic terrain structure for the user.

% The label map is a symbolic representation of the terrain, where each pixel of the map (each colour) represents a class. We can use this information to keep a sparse representation of the environment using medial axis transform shape descriptor, for which we will see how to treat in the next chapter about \glosses{EnvObj}.

% The training of the model can be done in the other direction, from height field to label map. \
% If we wish to extract its content into \glosses{EnvObj}, we can use the label map, skeletonise each class blob through medial axes, then attribute to each resulting polyline a class, that we will call, in the next chapter, \glosses{EnvObj}.

\paragraph{Other deep learning architectures}
We note that the the pix2pix original paper is more than a decade old, thus a comparison with modern models is necessary. However, the interactive condition (results in less than one second) eliminates most modern models. Secondly, the pix2pix model VRAM footprint is light ($<$4GB during training), while most models exceeds the capacity of low-end PCs (e.g.: pix2pixHD for $512\times512$ images is not usable on 4GB GPUs). Finally, we require the model to be flexible on the user inputs, meaning that the label map colors at any pixel may not exactly match a semantic map (see \cref{fig:coral-island-results-fuzzy}). For this last reason, models such as GauGAN requiring strict semantic inputs are unfit.

For illustration purpose, we proposed a simple prompt to three diffusion models inside the authors proposed interactive tools: OpenAI (GPT 5.5), Anthropic (Mistral Medium 3.5), and Flux2 (using the HuggingFace interface). We did not work on the prompt writting, it is possible that better results could be achieved with more work into it. Inference times using the proprietarian servers exceed greatly the interactive time limit ($>$6s for Flux2, $>$30s for Mistral Medium 3.5, and $>$45s for GPT 5.5). Results are shown in \cref{fig:training-llm-models}. 

\begin{figure}
    \autofitgraphics[]{LLM-to-terrain-trainings/input1.png, LLM-to-terrain-trainings/chatGPT/chatgpt-output1.png, LLM-to-terrain-trainings/mistral/mistral-output1.png, LLM-to-terrain-trainings/flux2/flux2-save2-output.png}
    \autofitgraphics[]{LLM-to-terrain-trainings/input2.png, LLM-to-terrain-trainings/chatGPT/chatgpt-output2.png, LLM-to-terrain-trainings/mistral/mistral-output2.png, LLM-to-terrain-trainings/flux2/flux2-save1-output.png}
    \autofitcaptions[]{Original input, GPT, Mistral, Flux2}
    \caption{Results of GPT5.5, Mistral, and Flux2. Prompt: "Transform the segmented image of this island into a grayscale heightmap of the island. Red is open ocean, dark blue is coral reef, light blue is lagoon, green is beach and yellow is the island center. The darkness of the image represents the age of the island, meaning that darker images creates more eroded and subsided islands compared to light images. The image must not be shaded. Each pixel of the output represent the height of the terrain."}
    \label{fig:training-llm-models}
\end{figure}